\section{Analytic deformation and pure Weil characters}
\label{sec:analytic-weil-sheaves}

This section excludes a coherent sheaf from the ordinary strongly
semiregular pure-character construction.  The distinction between a
coherent sheaf, an analytic perfect object and a global complex of
holomorphic vector bundles matters for this argument.
Bhattacharjee--Bhattacharya \cite[Theorem~A7]{BBK3Sixfolds} and
\cite[Theorem~10]{BBSecantObjects} state the exclusion for every
perfect complex.  Both proofs pass from formal deformations to
a global bounded complex of holomorphic vector bundles on nearby
tori.  That general passage fails by
Proposition~\ref{prop:analytic-strictification} below.
The coherent-sheaf assertion follows from the published analytic
deformation theorem of Buchweitz--Flenner.  We prove that assertion
and specify the additional implication needed for arbitrary complexes.
The Hodge-class and coherent-sheaf statements for very general
Weil tori are also proved in \cite[Lemmas~2--3 and Proposition~2]{BBSecantObjects},
following Voisin.  We include the vector-bundle conclusion and
the local analytic family description needed for the counterexample.

\begin{lemma}[Very general unpolarized Weil tori]
\label{lem:generic-weil-torus}
Let $A$ be a Weil abelian $2n$-fold for an imaginary quadratic
field $K$, with $n\ge2$.  Every neighborhood of $A$ in the local
space of balanced $K$-linear complex tori contains a torus $T$
with the following properties:
\begin{enumerate}[label=\textnormal{(\roman*)}]
\item its rational Hodge classes are generated by the rank,
the Weil plane $W_K$ in degree $2n$, and the top class;
\item every holomorphic vector bundle on $T$ has zero positive
rational Chern classes;
\item every coherent sheaf on $T$ has Chern character in
$H^0(T,\Q)\oplus H^{4n}(T,\Q)$.
\end{enumerate}
\end{lemma}

\begin{proof}
Put $V=H^1(A,\Q)$ and write
$V_\C=V_+\oplus V_-$ for the two $K$-eigenspaces, each of
dimension $2n$.  A balanced complex structure is given by
complementary $n$-planes $X,Y\subset V_+$, with
\[
 H^{1,0}=X\oplus\overline Y,
 \qquad H^{0,1}=Y\oplus\overline X.
\]
The deformation space is the open subset of
$\operatorname{Gr}(n,V_+)\times\operatorname{Gr}(n,V_-)$
consisting of $(X,Z)$ with $X\oplus\overline Z=V_+$.
It has complex dimension $2n^2$ and gives a smooth local analytic
family of marked tori.  As a real homogeneous space it is
$\operatorname{GL}_{2n}(\C)/
(\operatorname{GL}_n(\C)\times\operatorname{GL}_n(\C))$;
the preceding Grassmannian description specifies its deformation
complex structure.

The Hodge circle on $V_+$ acts by $e^{it}$ on $X$ and $e^{-it}$
on $Y$.  Its generator is nonzero and traceless.  As the
decomposition varies, the real Lie algebra generated by these
conjugate generators is a nonzero ideal in
$\mathfrak{sl}_{2n}(\C)$, regarded as a simple real Lie algebra.
It is therefore that whole Lie algebra.  Its complexification
acts on $V_\C$ as
$\operatorname{SL}(V_+)\times\operatorname{SL}(V_-)$.
In $\bigwedge^a V_+\otimes\bigwedge^b V_-$ an invariant occurs
only when each of $a,b$ is $0$ or $2n$.  The classes Hodge
throughout the parameter space are consequently exactly those
in (i).  For each other rational class its Hodge locus is a
proper analytic subset.  The countable union of these loci
cannot contain an analytic open set, by Baire's theorem.
This proves (i) on a very general torus in every neighborhood.

Choose a positive constant Hermitian metric making $X$ and
$\overline Y$ orthogonal.  Its Kähler class $\kappa$ belongs
to $V_+\wedge V_-$.  If $\alpha_\pm$ generates
$\bigwedge^{2n}V_\pm$, then
$\alpha_\pm\wedge\kappa=0$, since every term adds a vector
to an already exhausted eigenspace.  Thus $W_K\wedge\kappa=0$.
We have $\operatorname{NS}(T)_\Q=0$.  Moreover
$\operatorname{Hdg}^4(T)$ is zero for $n\ge3$, and for $n=2$
it is $W_K$ and is perpendicular to $\kappa^{2n-2}$.

Voisin's Proposition~A.3 and Corollary~A.4
\cite[pp.~1072--1074]{VoisinKahler} apply to exactly these two
conditions.  They say that every holomorphic vector bundle has
zero positive rational Chern classes, and that a torsion-free
coherent sheaf embeds in such a vector bundle with torsion
cokernel.  The input is the stable reflexive equality case of
Bando--Siu \cite[Theorem~3 and Corollary~3, p.~40]{BandoSiu}.
This proves (ii).

There is no proper positive-dimensional analytic subset of $T$.
Its fundamental class would be a nonzero rational Hodge class
of codimension $r$ with positive Kähler volume.  Statement (i)
leaves only $r=n$, but then its class is in $W_K$ and has
zero pairing with $\kappa^n$, a contradiction.  Every torsion
sheaf is therefore supported on finitely many points.  Remove
the torsion of an arbitrary coherent sheaf and apply Voisin's
embedding to the torsion-free quotient.  The embedding bundle
has character only in degree zero, and both torsion terms have
character only in top degree.  Additivity proves (iii).
\end{proof}

\begin{proposition}[Coherent pure-character exclusion]
\label{prop:coherent-pure-weil}
Let $A$ be a Weil abelian $2n$-fold, $n\ge2$, and let
$0\ne w\in W_K(A)$ be rational.  If a coherent sheaf
$\mathcal F$ satisfies
\[
 \operatorname{ch}(\mathcal F)-cw
   \in H^0(A,\Q)\oplus H^{4n}(A,\Q),\qquad c\ne0,
\]
then its full semiregularity map is not injective.  The same
assertion holds for a shift of such a sheaf.
\end{proposition}

\begin{proof}
Suppose the map were injective.  All components of the displayed
Chern character have horizontal Hodge extensions on the smooth
unpolarized Weil family in Lemma~\ref{lem:generic-weil-torus}.
Apply Buchweitz--Flenner \cite[Theorem~5.1, pp.~174--175]{BF},
with all Chern-character indices.  The central fibre is a compact
complex algebraic manifold, the base is smooth, and injectivity
of the full map is their required semiregularity hypothesis.

The conclusion needed here is the coherent sheaf constructed
in their proof on p.~179.  Proposition~5.9 and Lemma~5.10
extend $\mathcal F$ successively over the infinitesimal
neighborhoods of the base point.  A relative versal coherent-sheaf
deformation then has a compatible formal section.  Artin
approximation supplies a convergent section, whose pullback is
a coherent sheaf flat over a neighborhood in the analytic
family.  Its Chern character is the horizontal extension of
that of $\mathcal F$.  On a very general nearby fibre, however,
Lemma~\ref{lem:generic-weil-torus}(iii) makes its middle character
zero, whereas it must equal the nonzero transport of $cw$.
This is a contradiction.  A shift preserves self-Ext and changes
the character and trace map only by a sign.
\end{proof}

\begin{proposition}[Failure of global strictification]
\label{prop:analytic-strictification}
Injective semiregularity and a horizontal Hodge Chern character
do not imply that nearby nonprojective tori carry a global
bounded complex of holomorphic vector bundles with that character,
even when the initial perfect object has no negative self-Ext.
\end{proposition}

\begin{proof}
Take a point sheaf $E=\mathcal O_a$ on a projective Weil
$g$-fold $A$, where $g=2n\ge4$.  The local Koszul resolution
gives $\operatorname{Ext}^k(E,E)=\bigwedge^kT_aA$ and no
negative self-Ext.  In local coordinates its differential is
$\sum_j z_j\iota_j$; differentiating gives the Atiyah
representative $\sum_j dz_j\iota_j$.  Its $r$th power divided
by $r!$ contracts an $r$-vector by the usual exterior evaluation,
up to a nonzero sign.

Consider the component
\[
 \sigma_E^{g-2}:\operatorname{Ext}^2(E,E)
       \longrightarrow H^g(A,\Omega_A^{g-2}).
\]
Serre duality identifies its dual with contraction by
$\operatorname{At}(E)^{g-2}/(g-2)!$ from
$H^0(A,\bigwedge^{g-2}T_A\otimes\omega_A)$ to
$\operatorname{Ext}^{g-2}(E,E\otimes\omega_A)$.
The Koszul calculation identifies this with evaluation at $a$.
The tangent and canonical bundles are trivial, so evaluation
is an isomorphism.  Thus this component, and the full
semiregularity map, are injective.

The character of $E$ is the top class $[a]$, which stays Hodge
on every deformation of the torus.  Initially $E$ has a global
bounded resolution by algebraic vector bundles, since $A$ is
smooth and projective.  Every neighborhood in the unpolarized
Weil family contains a torus from
Lemma~\ref{lem:generic-weil-torus}(ii).  Every global bounded
complex of holomorphic vector bundles there has character
equal to the alternating sum of their ranks, hence has no
nonzero top component.  It cannot represent the transport of
$[a]$.  This proves the assertion.
\end{proof}

A point section still gives a coherent analytic point sheaf on
every nearby fibre.  It is an analytic perfect object, because
finite locally free resolutions exist locally on a smooth
complex manifold.  Proposition~\ref{prop:analytic-strictification}
concerns a single global resolution, not this local property.

For an arbitrary perfect complex, an analytic deformation theorem
producing analytic perfect objects on the varying family would
suffice for the argument of
Proposition~\ref{prop:coherent-pure-weil}.  Indeed its coherent
cohomology sheaves have zero intermediate Chern characters on
the very general fibre, and their alternating sum is the
character of the complex.  The global-resolution assertion
does not establish that implication.  We do not use the
arbitrary-perfect exclusion as a theorem, or apply the ordinary
coherent-sheaf exclusion to an invariant semiregularity map.
